% Figure -- Block C, step-invariant addressing. Column width (3.35in = 8.5cm),
% drawn at final size; never \resizebox'ed.
%
% Two granularities, top to bottom: what is built ONCE PER STEP (the
% compaction, its runs, the one block table) and what changes PER LAYER (only
% the base pointer of the slot handed to attention).
\begin{figure}[t]
\centering
\begin{tikzpicture}[x=1cm,y=1cm]
\tikzset{
  num/.style={circle, fill=black, text=white, inner sep=0.4pt,
              minimum size=3.6mm, font=\sffamily\bfseries\scriptsize},
  rc/.style={cell, minimum width=4.6mm, minimum height=4.2mm, fill=ddrF,
             draw=ddrD, font=\sffamily\scriptsize},
  ec/.style={cell, minimum width=4.6mm, minimum height=4.2mm, fill=white,
             draw=black!25},
  rl/.style={font=\sffamily\scriptsize\bfseries, text=ddrD},
}

%% ================= once per step ====================================
\fill[panel] (0,0.05) rectangle (8.40,4.55);
\node[hd, anchor=north west] at (0.08,4.50) {Once per step};
\node[sub, anchor=north east] at (8.33,4.50) {host, before layer 0};

% host DDR: three rider extents, scattered in the host pool
\node[font=\sffamily\scriptsize\bfseries, text=ddrD, anchor=west] at (0.10,3.62) {host DDR};
\node[sub, anchor=west] at (0.10,3.30) {block id};
\foreach \x/\id/\r in {2.00/7/A, 2.46/8/A, 3.82/23/B, 5.18/40/C, 5.64/41/C, 6.10/42/C}
  \node[rc] at (\x,3.55) {\id};
\foreach \x in {1.54,2.92,3.36,4.28,4.72,6.56}
  \node[ec] at (\x,3.55) {};
\node[font=\scriptsize, text=inkS] at (7.05,3.55) {$\cdots$};
\node[rl] at (2.23,3.98) {A};
\node[rl] at (3.82,3.98) {B};
\node[rl] at (5.64,3.98) {C};

% the compacted range [0,W)
\node[font=\sffamily\scriptsize\bfseries, text=hbmD, anchor=west] at (0.10,1.95) {slot layout};
\node[sub, anchor=west] at (0.10,1.63) {dense $[0,W)$};
\foreach \i/\id in {0/0,1/1,2/2,3/3,4/4,5/5}
  \node[rc] (c\i) at (3.08+\i*0.46,1.88) {\id};

% three runs = three DMA copies (H2D), the DMA engine does the compaction
\draw[h2d] (2.23,3.31) -- (3.31,2.12);
\draw[h2d] (3.82,3.31) -- (4.00,2.12);
\draw[h2d] (5.64,3.31) -- (4.92,2.12);
\node[font=\sffamily\scriptsize, text=h2dC, anchor=west, align=left] at (6.15,2.62)
  {3 runs\\$\Rightarrow$ 3 copies,\\not 6};
\draw[decorate, decoration={brace, amplitude=1.4mm, mirror}, black!55]
  (2.85,1.62) -- (3.77,1.62) node[midway, below=1.6mm, sub]{A};
\draw[decorate, decoration={brace, amplitude=1.4mm, mirror}, black!55]
  (3.79,1.62) -- (4.21,1.62) node[midway, below=1.6mm, sub]{B};
\draw[decorate, decoration={brace, amplitude=1.4mm, mirror}, black!55]
  (4.23,1.62) -- (5.61,1.62) node[midway, below=1.6mm, sub]{C};
\node[num] at (0.25,2.62) {1};
\node[sub, anchor=west] at (0.45,2.62) {sort + compact};

% the one block table
\node[ctl, fill=white, anchor=west, align=left, inner sep=1.4mm] (bt) at (0.10,0.58)
  {\textbf{one block table}\quad A\,[0,1]\;\; B\,[2]\;\; C\,[3,4,5]};
\node[num] at (bt.north west) {2};
\node[sub, anchor=west, align=left] at (5.98,0.58) {uploaded once,\\one slice per rider};

%% ================= per layer ========================================
\fill[panel] (0,-2.75) rectangle (8.40,-0.15);
\node[hd, anchor=north west] at (0.08,-0.20) {Per layer $\ell$ ($\times L$)};
\node[sub, anchor=north east] at (8.33,-0.20) {GPU};

% the k slots in HBM, each laid out as [0,W)
\node[font=\sffamily\scriptsize\bfseries, text=hbmD, anchor=west] at (0.10,-1.05) {$k$ slots};
\node[sub, anchor=west] at (0.10,-1.37) {in HBM};
\foreach \s in {0,1,2,3} {
  \node[hbm, fill=white, minimum width=1.30cm, minimum height=0.62cm] (s\s) at (2.05+\s*1.52,-1.20) {};
  \foreach \j in {0,...,5}
    \node[cell, minimum width=1.6mm, minimum height=3.2mm, fill=ddrF, draw=ddrD]
      at (1.65+\s*1.52+\j*0.16,-1.20) {};
  \node[font=\sffamily\scriptsize] at (2.05+\s*1.52,-0.74) {$s_{\s}$};
}
% base pointer moves l -> l+1
\draw[->, draw=ctlD, line width=0.9pt] (3.57,-2.15) -- (3.57,-1.55);
\draw[->, draw=ctlD, line width=0.9pt, densely dashed] (5.09,-2.15) -- (5.09,-1.55);
\node[font=\sffamily\scriptsize, text=ctlD, anchor=north] at (3.57,-2.13) {layer $\ell$};
\node[font=\sffamily\scriptsize, text=ctlD, anchor=north] at (5.09,-2.13) {layer $\ell{+}1$};
\node[num] at (0.25,-2.12) {3};
\node[sub, anchor=west] at (0.45,-2.12) {base $= s_{\ell \bmod k}$};
\node[sub, anchor=west, align=left] at (5.75,-2.20) {only the base pointer\\moves; same table};

%% ================= contrast ========================================
\node[bad, anchor=north west, align=left, text width=3.75cm, inner sep=1.3mm]
  (nv) at (0.0,-2.95)
  {\textbf{naive}: per rider \emph{block}\\$O(r\cdot S/b_s)$ host work, grows\\with context};
\node[cmp, anchor=north east, align=left, text width=3.75cm, inner sep=1.3mm]
  (us) at (8.40,-2.95)
  {\textbf{ours}: one slice per \emph{rider}\\$O(r)$ host work,\\independent of context};
\end{tikzpicture}
\caption{\textbf{Step-invariant addressing: addresses are fixed per step, contents
change per layer.} \stepn{1}~Once per step the riders' host extents are sorted and
compacted into a dense range; maximal contiguous runs become one DMA copy each (block-major
storage with $K$ and $V$ split keeps each run contiguous), so the copy engine
performs the compaction. \stepn{2}~The resulting block table is
built and uploaded once and serves all $L$ layers. \stepn{3}~Every slot has the same
layout, so advancing from layer $\ell$ to $\ell{+}1$ moves only the base pointer
handed to attention. Building the table one slice per rider cut its host cost from
0.275 to 0.0031\,ms per rider (ctx 512, A100), identical at ctx 2048.}
\label{fig:addressing}
\end{figure}
